% GENERATED FILE. Edit canonical lessons, then run npm run generate:books.
% canonical-source-hash: fnv1a64:b31d22fe8b11bccf
% canonical-lessons: HI-C398-musafir, HI-C398-kuli, HI-C398-dhhona, HI-C398-ladna, HI-C398-bithana

\chapter{Passengers and Porters}
\label{ch:hi-passengers-and-porters}
\hlchaptermodality{398}

\begin{chapteropening}
\textbf{By the end of this chapter:} \emph{I can talk about passengers and porters, the loads they carry and a seat found for someone.}

Complete the last lesson of chapter 398: I can talk about passengers and porters, the loads they carry and a seat found for someone.
\end{chapteropening}

\section[\texorpdfstring{\hi{मुसाफ़िर} (musāfir)}{मुसाफ़िर (musāfir)}]{\hi{मुसाफ़िर} (musāfir) — a traveller; a passenger}

\label{lesson:HI-C398-musafir}



\begin{quote}
Before the new one: say the Hindi for turning over (in bed), then the Hindi for a stretch (of the body).
\end{quote}

\subsection*{You'll want to know: \hi{मुसाफ़िर}}
\textbf{\hi{मुसाफ़िर}} — \emph{musāfir} — \textquotedblleft{}a traveller; a passenger\textquotedblright{}.

\begin{quote}
\hi{ट्रेन} \hi{में} \hi{बहुत} \hi{मुसाफ़िर} \hi{थे।}
\end{quote}

\begin{quote}
\emph{ṭren meṁ bahut musāfir the.}
\end{quote}

\begin{quote}
There were a lot of passengers on the train.
\end{quote}

Masculine or feminine, by the person meant.

\begin{cousinweb}
\textbf{\hi{मुसाफ़िर}} (\emph{musāfir}) is usually traced to Arabic \emph{musāfir}, \textquotedblleft{}one who travels\textquotedblright{}, which reached Hindi through Persian; on a train or a bus it is the everyday word for a passenger.
\end{cousinweb}

\begin{grammarlens}[title={What you've built}]
One more word for getting around.
\end{grammarlens}

\subsection*{Your turn}
\begin{itemize}
\raggedright
  \item \emph{Say it:} \emph{musāfir}
  \item \emph{Say it:} \emph{musāfir}, once more
  \item \emph{Say it:} \emph{ṭren meṁ bahut musāfir the.} — There were a lot of passengers on the train.
  \item \emph{Recall:} say the Hindi for turning over (in bed), then the Hindi for a stretch (of the body), then say \emph{musāfir} again
\end{itemize}

\begin{tcolorbox}[breakable,colback=teal!4,colframe=teal!35!black,title={Before you move on}]
What does \hi{मुसाफ़िर} mean? (\textquotedblleft{}A traveller; a passenger\textquotedblright{}.) Say it once more.
\end{tcolorbox}
\section[\texorpdfstring{\hi{कुली} (kulī)}{कुली (kulī)}]{\hi{कुली} (kulī) — a porter}

\label{lesson:HI-C398-kuli}



\begin{quote}
Before the new one: say the Hindi for a stretch (of the body), then the Hindi for a traveller.
\end{quote}

\subsection*{You'll want to know: \hi{कुली}}
\textbf{\hi{कुली}} — \emph{kulī} — \textquotedblleft{}a porter\textquotedblright{}.

\begin{quote}
\hi{कुली} \hi{ने} \hi{मेरा} \hi{सामान} \hi{उठाया।}
\end{quote}

\begin{quote}
\emph{kulī ne merā sāmān uṭhāyā.}
\end{quote}

\begin{quote}
The porter picked up my luggage.
\end{quote}

Masculine.

Station porters wear a red shirt and a numbered brass badge, so you can find the same one again.

\begin{cousinweb}
\textbf{\hi{कुली}} (\emph{kulī}) is often traced to Tamil \emph{kūli}, \textquotedblleft{}wages, hire\textquotedblright{} (DEDR 1905): the porter is the man you pay by the load. Other sources have been proposed, so the link is probable rather than certain.
\end{cousinweb}

\begin{grammarlens}[title={What you've built}]
One more word for getting around.
\end{grammarlens}

\subsection*{Your turn}
\begin{itemize}
\raggedright
  \item \emph{Say it:} \emph{kulī}
  \item \emph{Say it:} \emph{kulī}, once more
  \item \emph{Say it:} \emph{kulī ne merā sāmān uṭhāyā.} — The porter picked up my luggage.
  \item \emph{Recall:} say the Hindi for a stretch (of the body), then the Hindi for a traveller, then say \emph{kulī} again
\end{itemize}

\begin{tcolorbox}[breakable,colback=teal!4,colframe=teal!35!black,title={Before you move on}]
What does \hi{कुली} mean? (\textquotedblleft{}A porter\textquotedblright{}.) Say it once more.
\end{tcolorbox}
\section[\texorpdfstring{\hi{ढोना} (ḍhonā)}{ढोना (ḍhonā)}]{\hi{ढोना} (ḍhonā) — to haul; to carry (a heavy load)}

\label{lesson:HI-C398-dhhona}



\begin{quote}
Before the new one: say the Hindi for a traveller, then the Hindi for a porter.
\end{quote}

\subsection*{You'll want to know: \hi{ढोना}}
\textbf{\hi{ढोना}} — \emph{ḍhonā} — \textquotedblleft{}to haul; to carry (a heavy load)\textquotedblright{}.

\begin{quote}
\hi{कुली} \hi{दिन} \hi{भर} \hi{भारी} \hi{बोरियाँ} \hi{ढोते} \hi{हैं।}
\end{quote}

\begin{quote}
\emph{kulī din bhar bhārī boriyā̃ ḍhote haiṁ.}
\end{quote}

\begin{quote}
Porters haul heavy sacks all day long.
\end{quote}

Past, with \hi{ने} (he or she …): \hi{उसने} \hi{सामान} \hi{ढोया}, \emph{usne sāmān ḍhoyā}.

\begin{cousinweb}
\textbf{\hi{ढोना}} (\emph{ḍhonā}) is a tadbhava of Sanskrit \emph{ḍhaukayati}, \textquotedblleft{}brings near\textquotedblright{} (CDIAL 5610), which lists Hindi \emph{ḍhonā}, \textquotedblleft{}to convey\textquotedblright{}. Bringing a thing near became carrying it, load after load.
\end{cousinweb}

\begin{grammarlens}[title={What you've built}]
One more word for getting around.
\end{grammarlens}

\subsection*{Your turn}
\begin{itemize}
\raggedright
  \item \emph{Say it:} \emph{ḍhonā}
  \item \emph{Say it:} \emph{ḍhonā}, once more
  \item \emph{Say it:} \emph{kulī din bhar bhārī boriyā̃ ḍhote haiṁ.} — Porters haul heavy sacks all day long.
  \item \emph{Recall:} say the Hindi for a traveller, then the Hindi for a porter, then say \emph{ḍhonā} again
\end{itemize}

\begin{tcolorbox}[breakable,colback=teal!4,colframe=teal!35!black,title={Before you move on}]
What does \hi{ढोना} mean? (\textquotedblleft{}To haul; to carry (a heavy load)\textquotedblright{}.) Say it once more.
\end{tcolorbox}
\section[\texorpdfstring{\hi{लादना} (lādnā)}{लादना (lādnā)}]{\hi{लादना} (lādnā) — to load (onto)}

\label{lesson:HI-C398-ladna}



\begin{quote}
Before the new one: say the Hindi for a porter, then the Hindi for to haul.
\end{quote}

\subsection*{You'll want to know: \hi{लादना}}
\textbf{\hi{लादना}} — \emph{lādnā} — \textquotedblleft{}to load (onto)\textquotedblright{}.

\begin{quote}
\hi{मज़दूरों} \hi{ने} \hi{ट्रक} \hi{पर} \hi{सामान} \hi{लादा।}
\end{quote}

\begin{quote}
\emph{mazdūroṁ ne ṭrak par sāmān lādā.}
\end{quote}

\begin{quote}
The labourers loaded the goods onto the lorry.
\end{quote}

Past, with \hi{ने} (he or she …): \hi{उसने} \hi{गाड़ी} \hi{पर} \hi{बोरियाँ} \hi{लादीं}, \emph{usne gāṛī par boriyā̃ lādīṁ}.

\begin{cousinweb}
\textbf{\hi{लादना}} (\emph{lādnā}) is a tadbhava of Sanskrit \emph{lardayati}, \textquotedblleft{}loads\textquotedblright{} (CDIAL 10966); the r has dropped out and the vowel has lengthened in its place.
\end{cousinweb}

\begin{grammarlens}[title={What you've built}]
One more word for getting around.
\end{grammarlens}

\subsection*{Your turn}
\begin{itemize}
\raggedright
  \item \emph{Say it:} \emph{lādnā}
  \item \emph{Say it:} \emph{lādnā}, once more
  \item \emph{Say it:} \emph{mazdūroṁ ne ṭrak par sāmān lādā.} — The labourers loaded the goods onto the lorry.
  \item \emph{Recall:} say the Hindi for a porter, then the Hindi for to haul, then say \emph{lādnā} again
\end{itemize}

\begin{tcolorbox}[breakable,colback=teal!4,colframe=teal!35!black,title={Before you move on}]
What does \hi{लादना} mean? (\textquotedblleft{}To load (onto)\textquotedblright{}.) Say it once more.
\end{tcolorbox}
\section[\texorpdfstring{\hi{बिठाना} (biṭhānā)}{बिठाना (biṭhānā)}]{\hi{बिठाना} (biṭhānā) — to seat (someone); to make (someone) sit}

\label{lesson:HI-C398-bithana}



\begin{quote}
Before the new one: say the Hindi for to haul, then the Hindi for to load (onto).
\end{quote}

\subsection*{You'll want to know: \hi{बिठाना}}
\textbf{\hi{बिठाना}} — \emph{biṭhānā} — \textquotedblleft{}to seat (someone); to make (someone) sit\textquotedblright{}.

\begin{quote}
\hi{कुली} \hi{ने} \hi{दादी} \hi{को} \hi{सीट} \hi{पर} \hi{बिठाया।}
\end{quote}

\begin{quote}
\emph{kulī ne dādī ko sīṭ par biṭhāyā.}
\end{quote}

\begin{quote}
The porter sat my grandmother down in a seat.
\end{quote}

Past, with \hi{ने} (he or she …): \hi{उसने} \hi{मेहमान} \hi{को} \hi{कुर्सी} \hi{पर} \hi{बिठाया}, \emph{usne mehmān ko kursī par biṭhāyā}.

\begin{cousinweb}
\textbf{\hi{बिठाना}} (\emph{biṭhānā}) is the causative partner of \hi{बैठना}, \textquotedblleft{}to sit\textquotedblright{}, which goes back to Sanskrit \emph{upaviśati}, \textquotedblleft{}sits down\textquotedblright{} (CDIAL 2245): you \hi{बैठना} yourself, and you \hi{बिठाना} someone else.
\end{cousinweb}

\begin{grammarlens}[title={What you've built}]
One more word for getting around.
\end{grammarlens}

\subsection*{Your turn}
\begin{itemize}
\raggedright
  \item \emph{Say it:} \emph{biṭhānā}
  \item \emph{Say it:} \emph{biṭhānā}, once more
  \item \emph{Say it:} \emph{kulī ne dādī ko sīṭ par biṭhāyā.} — The porter sat my grandmother down in a seat.
  \item \emph{Recall:} say the Hindi for to haul, then the Hindi for to load (onto), then say \emph{biṭhānā} again
\end{itemize}

\begin{tcolorbox}[breakable,colback=teal!4,colframe=teal!35!black,title={Before you move on}]
What does \hi{बिठाना} mean? (\textquotedblleft{}To seat (someone); to make (someone) sit\textquotedblright{}.) Say it once more.
\end{tcolorbox}
